% !TEX root = ../main.tex

\providecommand{\kBKM}{\kappa_{\mathrm{BKM}}}
\providecommand{\kch}{\kappa_{\mathrm{ch}}}
\providecommand{\kcat}{\kappa_{\mathrm{cat}}}
\providecommand{\kfib}{\kappa_{\mathrm{fiber}}}
\providecommand{\DBKM}{\Delta^{\mathrm{BKM}}}
\providecommand{\PhiPhys}{\Phi^{\mathrm{phys}}}
\providecommand{\Z}{\mathbb{Z}}
\providecommand{\AdS}{\mathrm{AdS}}
\providecommand{\Tr}{\mathrm{Tr}}

\section*{Black-hole indices and CHL Borcherds weights}
\label{wn:sec:black-hole-indices-chl-borcherds}

\paragraph{CHL datum.}
Let
\[
  X^{(N)}=(K3\times E)/\Z_N,\qquad N\in\{1,2,3,4,6\},
\]
with $g_N$ a symplectic K3 automorphism and the elliptic factor
providing the free diagonal action. For charges $(Q,P)$ in the CHL
Narain lattice set
\[
  \Delta_N(Q,P)=Q^2P^2-(Q\cdot P)^2 .
\]
The protected black-hole index is scalar; the BKM denominator and any
Hall--Drinfeld construction are separate pieces of structure.

\paragraph{Theorem: scalar dyon coefficient.}
The 1/4-BPS indexed degeneracy is extracted from the physical Siegel
form by
\[
d_N(Q,P)=
(-1)^{Q\cdot P+1}
\oint_{\mathcal C_N}
\frac{
e^{-\pi i(Q^2 z_1+2(Q\cdot P)z_2+P^2 z_3)}
}{
\PhiPhys_N(Z)
}
d^3Z ,
\qquad
Z=\begin{pmatrix}z_1&z_2\\ z_2&z_3\end{pmatrix}.
\label{wn:eq:black-hole-scalar-dyon-coefficient}
\]
At $N=1$, $\PhiPhys_1=\Phi_{10}^{\mathrm{un}}=\Delta_5^2$ in the
unnormalised convention, so the scalar contour uses
$\Phi_{10}^{-1}=\Delta_5^{-2}$. In the Oberdieck--Pandharipande
normalisation one writes $D_5=64^{-1}\Delta_5$ and
$\Phi_{10}^{\mathrm{OP}}=D_5^2$, giving the scalar
$-(\Phi_{10}^{\mathrm{OP}})^{-1}=-4096\,\Delta_5^{-2}$ for the reduced
DT series. The scalar equality is a character identity. Compact Hall
correspondences and BKM root data require separate construction.

\paragraph{Theorem: leading black-hole entropy.}
The large-charge attractor asymptotic has leading term
\[
  \log d_N(Q,P)
  =
  \pi\sqrt{\Delta_N(Q,P)/N}+O(\log \Delta_N).
\label{wn:eq:black-hole-leading-entropy}
\]
This is the indexed form of the Bekenstein--Hawking entropy for the
CHL dyon chamber, with the $1/N$ determined by the CHL lattice and
orbifold normalisation. It is independent of any proposed Hall
realisation of the BKM denominator.

\paragraph{BKM logarithmic weight.}
Let $\DBKM_N$ be the primitive Borcherds denominator in the Vol III CHL
normalisation. Then
\[
  (c_1(0),c_2(0),c_3(0),c_4(0),c_6(0))=(10,8,6,4,2),
  \qquad
  \kBKM(\DBKM_N)=c_N(0)/2=(5,4,3,2,1).
\label{wn:eq:black-hole-bkm-constants}
\]
The BKM square-root normalisation contributes the automorphic weight
term
\[
  \kBKM(\DBKM_N)-\frac32=\frac{c_N(0)}2-\frac32
\label{wn:eq:black-hole-bkm-log-weight}
\]
before conversion to the physical contour convention. A comparison with
Sen's quantum entropy function must keep the square-root passage, the
doubled physical weight, the Gaussian measure, and the massless
zero-mode subtraction in the same convention. The unqualified
replacement of the physical denominator by the primitive BKM denominator
is an additional comparison problem.

\paragraph{Scalar partition function and DT character.}
At $N=1$ the reduced $K3\times E$ scalar satisfies
\[
  Z_{\mathrm{DT}}^{\mathrm{red}}(K3\times E)
  =
  -(\Phi_{10}^{\mathrm{OP}})^{-1}
  =
  -D_5^{-2}
  =
  -4096\,\Delta_5^{-2}.
\label{wn:eq:black-hole-op-scalar}
\]
This formula is a scalar partition-function theorem in the reduced DT
normalisation. Reading it as a graded Hall character is conditional on
constructing the oriented compact critical CoHA, the factorisation
pushforward, and the character map. The equality cannot supply the
primitive brackets, Serre relations, coproduct, or completion.

\paragraph{Integrality.}
The coefficients in \eqref{wn:eq:black-hole-scalar-dyon-coefficient}
are integer BPS indices after the contour and chamber are fixed.
At $N=1$ this is integrality of the $\Phi_{10}^{-1}$ coefficient
extraction. The square identity $\Phi_{10}=\Delta_5^2$ explains why
the same scalar coefficients are compatible with the primitive
denominator $\Delta_5$. A factorisation of the BPS index into primitive
BKM counts would require an independent primitive-counting construction.

\paragraph{Holographic heuristic.}
The AdS$_3$/CFT$_2$ and AdS$_2$ quantum-entropy descriptions give a
physics interpretation of the same protected index. A compact BPS
quantum group requires a construction beyond the scalar partition
function. A candidate Virasoro or Schur-sector action on
$\fg_{\Delta_5}$ must be constructed separately and compared with the
$E_1$ chiral algebra $\mathcal A_{X^{(1)}}$ at the level of stress
tensors. Negative central-charge numerology such as $c=-214$ is
conditional data for a non-unitary chiral algebra; a physical
Ryu--Takayanagi theorem for the D1--D5 CFT requires separate bulk
geometry.

\paragraph{Four invariants.}
For $N=1$ the $K3\times E$ constructions give
\[
  \{\kcat,\ \kappa_{\mathrm{ch}}^{\mathrm{Heis}},\ \kBKM,\ \kfib\}
  =
  \{0,3,5,24\}.
\]
The entries are independent: $\kcat(K3\times E)=0$ by Kunneth on
$\chi(\mathcal O_{K3})\chi(\mathcal O_E)=2\cdot 0$;
$\kappa_{\mathrm{ch}}^{\mathrm{Heis}}=3$ belongs to the
Heisenberg--Mukai specialisation; $\kBKM(\Delta_5)=5=c_1(0)/2$ belongs
to the primitive Borcherds denominator; and $\kfib=24$ is the Mukai
  lattice rank of the K3 fibre. The compact CY$_3$ output has
$\kch(\mathcal A_{X^{(1)}})=0$ by the odd-dimensional Hodge
supertrace. The additive formula
$\kBKM=\kch+\chi(\mathcal O_{\mathrm{fiber}})$ gives $5=0$ at $N=1$.

\paragraph{Operadic scope.}
Since $\dim_{\mathbb C}X^{(N)}=3$, the specialised chiral algebra
$\mathcal A_{X^{(N)}}$ is native $E_1$. An $E_2$ structure can appear
only after passing to a derived chiral centre, a Drinfeld double, or a
completed Hall construction with the required continuity and pairing
data.

\paragraph{Proof obligations.}
The following claims cannot be used as theorem-level input without
additional construction: the compact Hall--Drinfeld realisation of
$\DBKM_N$; the equality between a reduced DT scalar and a compact Hall
character beyond the conditional character-map statement; the
conversion between \eqref{wn:eq:black-hole-bkm-log-weight} and the full
physical logarithmic correction in one convention; and any direct
class-$\mathcal S$, CoHA, or higher-spin theorem derived from
$\Phi_{10}^{-1}$ alone.

\paragraph{Sources.}
Dabholkar--Harvey 1989; Strominger--Vafa 1996; Dijkgraaf--Verlinde--
Verlinde 1997; Dijkgraaf--Moore--Verlinde--Verlinde 1997; Maldacena--
Moore--Strominger 1999; Jatkar--Sen 2006; Sen 2007, 2008, and 2011;
Bekenstein 1973; Hawking 1975; Gritsenko--Nikulin 1995; Gritsenko
1999; Borcherds 1998; Oberdieck--Pandharipande 2016.
